\chapter{질문과 답}
\newcommand{\QA}[2]{\begin{qa}{Q. #1}#2\end{qa}}

팀 스터디(2026년 11월 16일)에서 보고서를 함께 읽으며 나온 질문과 답이다. 발표와 질의응답 대비용이다. 숫자는 모두 train(2013--16년 36개월 대출) 결과다.

\section{전체 구조}
\QA{결국 우리 과정은 어떻게 흘러가나?}{
\begin{enumerate}[leftmargin=*]
\item 대출마다 부도확률 $p$를 구한다. 모델은 그 시점에 알 수 있던 과거 대출로만 학습한다.
\item 점수로 줄을 세운다. 기대수익 순(V1)은 $s=(1-p)r-p\cdot\mathrm{LGD}-tb$, 부도확률 순(V2)은 $p$가 낮은 순이다.
\item 발행월마다 그 달 대출 중 상위 $q$를 승인한다.
\item 그 달 승인 대출의 현금흐름을 합쳐 발행월 묶음 IRR을 하나 구한다(48개월이면 48개).
\item 매달 초과수익 $=$ (묶음 IRR $-$ 같은 달 같은 만기 국채) \ensuremath{\times} 승인 금액 비중.
\item 48개 초과수익의 평균 \ensuremath{\div} 표준편차 $=$ 포트폴리오 Sharpe.
\item 같은 비율 무작위 선택과 비교해 고르는 실력이 있는지 본다.
\end{enumerate}}

\QA{IRR을 Sharpe의 분자에 넣는 것인가? ``IRR $-$ $s$''를 계산하나?}{IRR 자체가 아니라 (IRR $-$ 국채) \ensuremath{\times} 비중이 분자에 들어간다. $s$는 승인할 때 쓰는 예상 점수(고르는 도구)이고, IRR은 끝난 뒤 실제 현금흐름으로 잰 수익률(채점 도구)이다. 둘을 빼거나 섞지 않는다. 예상 점수로 학생을 뽑고 실제 점수로 그 선택을 채점하는 구조다.}

\QA{기대수익 점수 $s$는 우리가 만든 식인가?}{아니다. 문헌의 기대수익 점수(expected return score)다. 안 망하면 받는 이자, 빼기 망하면 잃는 원금, 빼기 국채 기회비용이다. 1년 단위 근사라 정확한 예상 수익이 아니라 순위용으로만 쓴다.}

\QA{IRR만 보면 되지 왜 Sharpe가 필요한가?}{평균 초과수익만 키우면 답이 ``전부 승인''(4.86\%p)이다. 48개월 내내 대출이 국채보다 나았으니 많이 살수록 평균이 커지고, 부도 위험은 전혀 보지 않게 된다. Sharpe는 흔들림(위험)을 분모에 넣어 준다. 반대로 Sharpe만 최대로 하면 ``가장 안전한 10\%''(3.82)가 답이라 돈의 90\%가 논다. 그래서 Sharpe를 주 지표로, 평균 초과수익을 두 번째 축으로 본다.}

\QA{우리 목표는 무엇인가? Sharpe 극대화인가 수익 극대화인가?}{권하는 해석은 \textbf{``Sharpe를 무작위와 전부 승인보다 확실히 높이면서 수익 규모를 크게 잃지 않는 것''}이다.
\begin{itemize}[leftmargin=*]
\item 평균 초과수익 극대화: 전부 승인. 선별이 필요 없고 위험을 보지 않는다.
\item Sharpe 엄격 최대화: 부도확률 순 10\%. 규모가 작고, 흔들림이 거의 0이라 숫자가 불안정하며, 18개 중 최대값이라 운이 섞였다.
\item Sharpe 개선 $+$ 규모 유지: 부도확률 순 80\%(3.06, 평균 4.30\%p). 무작위(2.72), 전부 승인(2.71)보다 높고 규모는 전부 승인의 약 90\%다.
\end{itemize}
test 판정 H10(같은 비율 무작위보다 Sharpe가 높은가)도 이 목표와 맞는다. 다만 80\%는 v2.2 규칙이 고른 \textbf{현재 후보}이며, 팀이 Sharpe 엄격 최대화로 정하면 test 전에 v2.3으로 다시 고정한다.}

\QA{LC에서 왜 ``포트폴리오''로 보나?}{LC는 플랫폼이다. 투자자는 대출 하나를 통째로 빌려주지 않고 여러 대출에 조각(25달러 단위)으로 나눠 투자한다. LC는 투자자 수수료 1\%를 뗀다. 그래서 대출 묶음의 성과가 투자자 입장의 성적이다.}

\section{요약 문장}
\QA{``매년 약 1\%p 더 번다''는 얼마인가?}{1\%p는 수익률 0.01이다. 연 5\%가 6\%가 되는 정도이고, 1억 원을 1년 굴리면 약 100만 원 차이다. 95\% 구간 0.95--1.11은 운을 감안해도 그 범위라는 뜻이다.}

\QA{롱숏은 등급마다 하나, 전체를 $p$로 줄 세우나?}{\textbf{등급마다 따로} 한다. 매달, 등급마다 기대수익 점수 상위 30\%의 IRR에서 하위 30\%의 IRR을 뺀다. 전체를 줄 세우면 위쪽이 거의 A등급이라 ``좋은 등급을 산 효과''가 섞인다. 등급 안에서 빼야 고르는 실력만 남는다(+2.38\%p, 48개월 전부 양수). LC 대출은 공매도가 안 되므로 실제 전략이 아니라 측정 도구다.}

\QA{회사채(BAA)는 애플 같은 안 망할 회사인가? 국채는 몇 년물인가?}{BAA는 Moody's 투자등급 중 가장 낮은 등급(Baa) 회사채 수익률 지수다. AAA보다 위험하고 금리가 높다. 그래서 ``투자자의 요구수익률이 높을 때''의 허들로 쓴다. 국채는 대출 만기에 맞춰 36개월 대출에 3년물(GS3)을 쓴다.}

\QA{``새 계좌를 적게 연''은 무슨 뜻인가?}{\ft{acc_open_past_24mths}: 신청자가 최근 24개월 동안 새로 연 신용 계좌(카드, 대출 등) 수다. 짧은 기간에 계좌를 많이 열면 돈이 급하다는 신호라 부도와 연결된다. 이 규칙(V3)은 데이터를 보고 찾은 사후 규칙이라 주 제안에서 빼고 test에서 확인할 도전자로만 둔다.}

\QA{등급은 $p$ 계산에만 쓰이나?}{세 곳에 쓰인다. (1) $p$ 모델의 입력(등급, 세부등급, 금리). (2) 롱숏과 V3의 등급 안 비교. (3) 성과 분해(Brinson). V1, V2의 승인 컷은 등급을 나누지 않고 발행월 전체에서 자른다.}

\section{수익률과 IRR}
\QA{2장 예시(1만 달러, 332.14 \ensuremath{\times} 10, 회수 500)는 누구 입장인가?}{투자자(빌려준 쪽)다. 0개월에 10,000을 내주고 1--10개월에 월 상환액 332.14를 받다가 15개월에 회수금 500을 받았다. 이 현금흐름의 IRR이 연 $-76.8\%$다.}

\QA{단순 연환산 $-27.4\%$의 36은 어디서 왔나?}{대출 만기 36개월이다. $(3821.43/10000)^{12/36}-1$로, 실제로는 15개월 만에 끝난 손실을 3년에 걸쳐 늘려서 작아 보인다. 지난 우승 논문의 계산 방식이며, 비판할 때는 논문 출처를 함께 적는다.}

\QA{우리 IRR은 코호트 IRR인가? 근거는?}{발행월 묶음(코호트) IRR이다. 같은 달 승인 대출의 현금흐름을 먼저 더한 뒤 IRR을 한 번 구한다. 근거는 다음과 같다.
\begin{itemize}[leftmargin=*]
\item 투자자 계좌가 실제로 번 수익률이다(금액 가중). 건별 IRR의 평균은 1,000달러와 30,000달러 대출을 똑같이 센다.
\item 건별 IRR은 부도 시 $-77\%$처럼 극단값이 커서 평균이 몇 건에 휘둘린다.
\item 돈이 들어온 시점(조기상환, 부도)을 반영한다.
\item 매달 승인한다는 결정 단위와 맞고, 48개 값이 생겨 흔들림을 잴 수 있다.
\end{itemize}}

\QA{IRR을 구할 때 금리가 필요한가? 명목인가 실질인가?}{IRR 자체에는 시장 금리가 필요 없다. 현금흐름만으로 현재가치 합이 0이 되는 할인율을 역산한다. 모든 값이 명목이고 허들(국채)도 명목이라 초과수익에서 인플레이션은 지워진다.}

\QA{조기상환은 손해인가? 금리가 오르면 오히려 이득 아닌가?}{맞는 지적이다. 조기상환의 손익은 재투자 금리에 달려 있다. 금리가 오르면 일찍 받은 돈을 더 높은 금리에 다시 넣을 수 있어 이득이다. IRR에는 같은 수익률로 재투자한다는 가정이 암묵적으로 들어 있어 이 효과를 따로 보지 않았다. 2013--16년은 저금리가 안정적이라 영향이 작았을 것으로 보지만 한계로 적었다(10장).}

\QA{현금흐름이 추정이라는 게 왜 한계인가? 월 상환액은 일정한데.}{월 상환액은 일정해도 \textbf{실제로 언제 냈는지}는 모른다. 데이터에는 받은 총액과 마지막 납부월만 있다. 연체 후 몰아서 낸 경우, 조기상환 시점, 부도 전 부분 납부는 추정이다. IRR은 시점에 민감하므로 한계로 적었다.}

\section{데이터와 표본}
\QA{결과 라벨(완제, 조기상환, 부도)은 어떻게 만들었나?}{원본 상태는 Fully Paid(완제)와 Charged Off(부도)다. 완제 가운데 만기 전에 끝난 것을 조기상환으로 구분해 분석에 썼다. 진행 중(Current, Late)은 결과를 모르니 학습과 채점에서 뺀다.}

\QA{종결 편향은 무엇인가?}{발행 후 1--3년밖에 안 된 대출 중 이미 끝난 것만 보면, 빨리 망한 대출은 이미 결과가 나와 있고 정상 상환자는 아직 갚는 중이다. 그래서 부도가 실제보다 많이 잡히고, ``부도만 피하는'' 규칙이 똑똑해 보인다. 그래서 만기 도달 비율이 낮은 연도는 따로 보고한다.}

\QA{생존분석은 왜 안 썼나?}{생존분석은 아직 안 끝난 대출을 버리지 않고 쓰는 정석 방법이다. 우리는 시점 학습과 종결 대출만 쓰는 방식을 택했고, 진행 중 대출을 활용한 실험은 효과가 없었다. 시간과 복잡도 때문에 쓰지 않았고, 한계 겸 추후 연구로 적었다(10장).}

\QA{MNAR, WOE는 무엇인가?}{MNAR(missing not at random)은 값이 비어 있는 것 자체가 정보라는 뜻이다. 예를 들어 종결 여부는 부도 여부와 관련이 있다. WOE(weight of evidence)는 변수의 각 구간이 정상과 부도를 얼마나 가르는지 로그 비율로 나타낸 값이다(4장).}

\section{부도확률 $p$}
\QA{FICO는 무엇이고 왜 690점인가?}{FICO는 미국 개인 신용점수(300--850)다. 신청 시점 값이라 누수가 아니다. 690은 설명용 예시이고, 실제 기준점은 모델이 데이터에서 찾는다. 나무가 수백 그루라 기준점도 나무마다 다르다.}

\QA{$p$는 왜 실제보다 낮게 나왔나?}{학습 시점(예: 2013년 1월 1일)에 쓸 수 있는 직전 해 발행분은 이미 끝난 대출뿐이다. 부도는 연체 후 확정까지 오래 걸리므로, 1년 안에 끝난 대출은 대부분 조기상환이다. 모델은 ``요즘 대출은 안전하다''고 배워 $p$가 절반쯤 낮게 나왔다(평균 6.4\% vs 실제 12.2\%). 일찍 하교한 학생만 보고 반 전체가 우등생이라고 판단한 셈이다.}

\QA{종결 편향(부도 과대)과 방향이 반대 아닌가?}{대출의 나이가 다를 뿐 둘 다 사실이다. 나이 1년 미만의 끝난 대출에는 조기상환이 많고(부도는 아직 확정 전), 1--3년의 끝난 대출에는 부도가 많다(정상 상환은 아직 만기 전). $p$를 낮춘 것은 앞의 것이다.}

\QA{$p$가 틀렸는데 왜 결과는 괜찮았나? 오차 상쇄는 구조적 설명인가?}{상쇄는 사후 관찰이지 보장된 구조가 아니다. 직접 근거는 두 가지다. (1) V2는 $p$의 순서만 쓰므로 크기가 같은 비율로 틀리면 결과가 그대로다. (2) V1은 금리가 함께 들어가 $p$의 크기가 순서를 일부 바꾸므로, $p$를 실제 수준으로 맞춰 다시 골랐다. 승인 대상의 88--92\%가 같았다. 순서까지 망가질 만큼 틀리면 문제다(CatBoost 2016년, 평균 3.6\% vs 실제 15.9\%).}

\QA{$p$를 왜 10분위로, 등급 안에서는 왜 셋으로 나눴나?}{10분위는 확률 신뢰도를 볼 때 쓰는 관례다(분위마다 약 5.2만 건). 등급 안 3분할은 등급마다 $p$ 하위, 중간, 상위 1/3이다. 분할 수는 임의이며, 결론은 ``같은 등급 안에서도 $p$가 높을수록 실제 부도율이 오른다''는 방향이다.}

\QA{AUC, seed, 부트스트랩은 무엇인가?}{AUC는 부도인 사람을 정상인 사람보다 위험하게 줄 세울 확률이다(0.5는 동전, 1은 완벽, 우리 0.63--0.68). seed는 무작위 계산의 출발값이고, 바꿔도 결과가 유지되는지 본다. 부트스트랩은 가진 표본을 다시 뽑아 ``운이 달랐다면''을 수천 번 흉내 내 구간을 만드는 방법이다. 우리는 6개월 덩어리로 다시 뽑아 이웃 달의 상관을 보존한다.}

\section{비교와 검증}
\QA{왜 ``같은 비율 무작위''와 비교하나? 무작위는 전체에서인가 등급 안에서인가?}{발행월 전체에서 같은 비율을 무작위로 고른다. 같은 달, 같은 비율이라 그달 시장 상황이 양쪽에 똑같이 들어가 지워지고, 차이에는 ``무엇을 골랐나''만 남는다.}

\QA{``정답을 섞은 가짜 모델은 0''은 무슨 뜻인가?}{부도 여부(정답)를 무작위로 섞어 학습시키면 모델은 아무것도 배울 수 없다. 그 모델로 고르면 선별 가치가 0 근처여야 하고, 실제로 그랬다. 평가 코드가 가짜 성과를 만들지 않는다는 점검이다.}

\QA{선별 가치의 ``9개 비율 평균''의 9는?}{승인 10\%, 20\%, \ldots, 90\%의 9개다. 유리한 비율 하나만 고르지 못하게 평균을 낸다.}

\section{Sharpe}
\QA{``건별''이란 무엇인가? 세 정의는 어떻게 다른가?}{건별은 대출 한 건을 데이터 한 줄로 본다는 뜻이다.
\begin{itemize}[leftmargin=*]
\item 건별, 거절 $=0$: 승인하지 않은 대출을 0으로 센다. 승인비율만 줄여도 값이 작아진다.
\item 건별, 승인만: 승인 대출만 본다. 0.10--0.34로, 조교가 말한 0.2--0.4와 크기가 맞는다.
\item 발행월 포트폴리오(주 지표): 한 달을 한 줄로 본다. 약 2--3.5.
\end{itemize}
어떤 정의로 채점될지는 따로 확인하지 않는다(평가가 과정과 논리 중심이므로). 정의를 왜 골랐는지 밝히고 건별 값을 참고로 함께 적는다.}

\QA{``흔들림도 0.3배''는 무슨 뜻인가? $\sqrt{0.3}$ 아닌가?}{모든 값에 0.3을 곱하면 표준편차도 정확히 0.3배다. 예: 10, 4, 7\% \ensuremath{\to} 3, 1.2, 2.1\%에서 오르내리는 폭이 6\%p \ensuremath{\to} 1.8\%p. 그래서 평균 \ensuremath{\div} 표준편차에서 0.3이 지워진다. $\sqrt{\ }$은 \textbf{대출 건수}가 바뀔 때의 효과다. 묶음에 대출이 적으면 개별 부도 운이 덜 희석되어 묶음 IRR이 더 흔들린다(대략 $1/\sqrt{\text{건수}}$). 그래서 승인 10\%의 Sharpe가 운에 더 민감하다(무작위 seed 간 표준편차 10\%: 0.23, 80\%: 0.02).}

\QA{포트폴리오 Sharpe가 왜 이렇게 크게(2--3.5) 나오나? 공격받지 않을까?}{세 가지 이유다. (1) 대출은 만기까지 보유하고 시장 가격으로 평가되지 않아, 주식처럼 가격이 오르내리는 흔들림이 없다. (2) 수백 건을 묶으면 개별 부도가 희석된다. (3) 2013--16년은 경기침체가 없었다. 그래서 주식 Sharpe(0.5 안팎)와 비교하면 안 되고, 우리 전략끼리 순위를 매기는 숫자로만 쓴다. 발표에서는 값 자체를 내세우지 않고 ``같은 비율 무작위보다 확실히 높다''까지만 주장한다.}

\QA{1년 단위로 묶으면 더 안정적이지 않나?}{1년 단위면 값이 4개(2013--16)뿐이라 표준편차를 믿을 수 없다. 또 묶는 기간이 길수록 흔들림이 작게 잡혀 위험이 과소평가된다. 월 단위는 실제 결정 단위(매달 승인)와도 맞는다.}

\QA{왜 연도 단위가 아니라 발행월마다 자르나?}{매달 승인 비중이 비슷해야 비중이 Sharpe에서 지워진다. 연도 단위로 자르면 고금리 대출이 몰린 달은 많이, 적은 달은 적게 승인되어 비중이 들쭉날쭉하고 Sharpe가 깎인다(2.46 vs 2.91). 12월 대출까지 본 뒤 1월 승인을 정하는 셈이라 실전에서도 불가능하다.}

\QA{Sortino는 무엇인가?}{Sharpe의 분모를 아래로 떨어진 흔들림만으로 바꾼 지표다. 건별로는 순위가 Sharpe와 같아 새 정보가 없었다. 포트폴리오로는 48개월 모두 국채보다 높아 아래쪽 흔들림이 0이라 계산되지 않는다.}

\section{팀 내 비판과 앞으로}
\QA{``좋은 시절만 봤다''는 한계는 어떻게 보완하나? 코로나 시기를 쓰면?}{데이터에 2017--20년 발행분이 있어 2017--19년 대출이 2020년 코로나 시기를 진행 중으로 겪었다. 다만 대부분 아직 진행 중이고, 데이터는 2020년 어느 시점의 상태이며, LC는 2020년에 투자자용 상품을 줄이고 닫았다. 그래서 Sharpe 비교가 아니라 \textbf{보조 스트레스 점검}으로만 권한다. 질문은 ``나쁜 시기에도 점수 상위 대출이 덜 연체·\allowbreak{}부도났나''이고, 판정 기준을 먼저 적는다.}

\QA{다른 팀원이 제안한 단순 규칙과 우리 규칙은 어떻게 다른가?}{그 규칙의 정확한 정의(부도확률 순서인지, 다른 변수인지)를 받아 V2, V3와 같은 표에서 비교해야 한다. 아직 하지 않았다.}

\QA{Brinson, 롱숏 같은 부가 분석은 본문에 다 넣어야 하나?}{큰 줄기만 본문에 두고, 이해가 덜 되거나 시간이 부족하면 부록이나 추후 연구로 돌린다. 롱숏은 핵심 숫자 하나(+2.38\%p, 48개월 전부 양수)만 본문에 남겨도 충분하다.}

\QA{특정 숫자를 집어서 발표해도 되나?}{된다. 다만 (1) train 숫자임을 표시하고, (2) 1등 숫자 하나만 앞세우지 말고 규칙과 함께 보이며, (3) 구간을 함께 쓴다. 예: ``Sharpe 3.06''보다 ``약 3, 무작위보다 +0.34(0.30--0.55)''.}
